%%
% BIThesis 研究生学位论文模板 The BIThesis Template for Graduate Thesis
% This file has no copyright assigned and is placed in the Public Domain.
%%

\chapter{绪论}
\label{chap:intro}

\section{研究背景、目的与意义}

在数字化与智能化应用逐步渗透专业场景的过程中，业务运行、信息传播等环节持续积累起大量知识数据。这些数据常常混含结构化、半结构化和非结构化内容，来源多、表达方式不统一，术语密集。由于数据形态与来源的差异，知识往往散落在不同平台、不同文档形态和不同表达层级中，难以被统一组织、持续维护和直接利用。知识图谱作为一种面向语义组织的知识表示方式，并不等同于事实三元组的简单堆叠。它整合了模式约束、实体标识、上下文语义与质量管理机制\cite{Hogan2021KnowledgeGraphs}。相比通用知识图谱，领域知识图谱更依赖于专业概念体系、本体约束与术语规范，其构建过程对知识来源的质量、结构的一致性和持续更新能力提出了更高要求\cite{AbuSalih2021DomainSpecificKGSurvey}。在多源异构文本场景下，如何将分散于各类文档中的原始知识转化为可计算、可复用、可追踪的结构化形式，是领域知识图谱构建必须首先面对的基础问题\cite{Sun2024Docs2KG,Zhong2023AutomaticKGConstructionSurvey}。

传统的“关键词检索-片段返回”方式在处理简单事实查询时效率尚可，但在高专业场景中面对复杂知识需求时，局限性明显。许多实际问题并非以显式关键词的形式出现，而是隐含在实体之间的关系链条、事件之间的因果规则，以及知识随时间的演化之中。片段式返回虽能提供局部文本证据，却难以支撑面向关系的知识组织、依据追溯和语义级检索。要真正提升领域知识的组织与利用能力，仅靠检索并不充分，还需要将文本中的实体、关系、事件、规则等要素显式抽取出来，组织成可查询、可维护、可持续更新的知识底座。

知识获取是原始文本进入知识图谱的入口\cite{Zhong2023AutomaticKGConstructionSurvey}。在获取阶段，知识抽取，尤其是实体识别与关系抽取，直接决定了图谱中节点、关系及语义链路能否被正确建立\cite{Zhao2024RelationExtractionSurvey}。基于规则模板、特征工程和监督学习的传统方法在特定任务中具有一定可解释性，但通常依赖大量人工规则或高质量标注数据。面对领域文本中新概念、新表达的持续出现，这类方法普遍面临开发成本高、维护周期长、跨场景迁移困难等问题\cite{Zhao2024RelationExtractionSurvey,Deng2024LowResourceIESurvey}。

大语言模型在复杂文本理解、少样本迁移和指令跟随方面展现出明显优势，为领域知识抽取提供了新的技术可能。不过，现有研究也表明，大语言模型在知识图谱构建中的价值更多体现在语义理解增强、流程加速和弱监督支持上，并不天然等同于稳定可靠的端到端知识抽取器\cite{Zhu2024LLMsForKGConstructionAndReasoning}。尤其是在高专业、高合规的场景中，对抽取方法的评价不能仅停留于局部精度，还必须同时考虑成本、吞吐、部署方式、可审计性，以及其与图谱模式的衔接能力。

因此，面向领域知识图谱构建的知识抽取研究，正在从“能否抽取”转向“如何实现高质量、可控且可部署的抽取”。本文讨论的“高质量”，不只是抽取结果在表面上的合理性，更强调训练数据与抽取结果在事实性、完整性、一致性、可利用性上能够满足后续建模需要。“可控”则不仅要求输出格式规整，还要求结果满足本体约束、具备结构合法性，并能够支持稳定入库与持续更新\cite{Tan2024LLMDataAnnotationSurvey,Klie2024AnalyzingAnnotationQualityManagement}。在长期运行与本地部署场景下，参数高效微调与小参数模型在资源占用、时延、维护成本等方面具有现实优势\cite{Hu2022LoRA,Dettmers2023QLoRA,Han2024PEFTSurvey,Lu2024SmallLanguageModels}。

基于上述分析，本文将研究对象限定为面向领域知识图谱构建的高质量可控知识抽取问题，具体围绕三个方面展开：一是 Teacher 端高质量银标数据的构建方式，二是 Student 端参数高效微调与本体约束下的可控抽取，三是知识图谱构建与下游应用场景的验证。理论层面，本文试图推动知识抽取研究从单纯追求局部精度，转向关注高质量银标数据、结构合法性和稳定入库等更贴近真实知识组织过程的问题。应用层面，本文为领域知识图谱的稳定构建、持续更新和工程化应用验证，提供一种相对可靠的结构化知识获取路径。需要说明的是，知识图谱问答系统并非本文的研究重点，下游应用仅作为验证所提方法工程价值的场景。

\section{国内外研究现状}

与本文相关的研究，大致可以沿着三条线索展开：领域知识图谱如何构建，知识抽取方法如何演进，以及面向真实场景的高质量可控知识抽取如何实现。这三者并非先后承接，而是相互嵌套。领域知识图谱规定了知识组织的基本框架与约束条件，知识抽取则是将文本等异构数据转化为结构化知识的核心手段，而高质量可控抽取更直接回应了实际部署中的稳定性、准确性与可维护性问题。

\subsection{领域知识图谱构建研究现状}

在领域知识图谱构建这一方向上，现有研究已将其与通用知识图谱做了明确区分。通用图谱追求广度与跨领域连接能力，而领域图谱更强调专业概念体系、本体约束与术语规范，其目标不是规模扩张，而是知识表示能否贴合特定任务的需求\cite{AbuSalih2021DomainSpecificKGSurvey}。一个领域图谱是否可用，往往不取决于它存了多少三元组，而在于它的实体划分是否契合业务逻辑、关系定义是否一致、后续能否持续维护和更新\cite{Hogan2021KnowledgeGraphs}。

从构建流程来看，知识获取、精炼与演化是基本环节\cite{Zhong2023AutomaticKGConstructionSurvey}。但对领域图谱而言，这几个环节必须与本体设计、元数据管理、质量控制等绑定在一起，否则即便建起来了，也很难在真实任务中稳定运行。已有研究反复指向同一个结论：领域知识图谱的难点不在存储层，而在于上游，如何从多源异构数据中稳定地识别出关键实体、关系和事件要素，并完成归一、融合与增量更新\cite{AbuSalih2021DomainSpecificKGSurvey,Zhong2023AutomaticKGConstructionSurvey}。换句话说，知识抽取的质量，直接决定了图谱能达到的上限。

这种难点的形成，与领域知识本身的特性密切相关。第一，数据来源高度异构：既有结构化字段，也有半结构化记录和非结构化文本，不同来源在命名方式、粒度划分、表达习惯上差异明显\cite{Sun2024Docs2KG}。第二，术语体系复杂，同一概念常存在多种表述，不同概念又可能在字面上相近但语义不同，这使得实体识别与消歧成本显著上升。第三，领域知识更新快，新概念、新事件、新规则不断出现，图谱必须具备持续的演化能力。第四，许多领域图谱服务于高专业、高合规场景，对关系定义、命名规范和入库规则的一致性要求远高于通用图谱。以上四个因素叠加在一起，使得“把数据变成图谱”这一过程远没有看上去那么直白。

目前，已有研究在多个高专业场景中开展了领域知识图谱构建的探索，表明自动化知识组织确实存在现实需求\cite{Sun2024Docs2KG,vanCauter2024OntologyGuidedKGC,Zhang2024EDC}。但若要进一步提升知识组织的质量，研究视角还需要向上游收束，回到多源异构文本场景中的知识抽取问题上，这恰恰是目前制约领域知识图谱走向深度应用的关键环节。

\subsection{知识抽取方法研究现状}

知识抽取是构建领域知识图谱的基础性工作，涉及实体识别、关系抽取以及实体-关系联合抽取等任务。从方法演进的角度看，这一领域大致经历了三个阶段。早期研究以规则和特征工程为主，代表性工作包括基于词汇-句法模式的关系模板、条件随机场、依存路径特征以及远监督方法\cite{Hearst1992Hyponyms,Lafferty2001CRF,Bunescu2005ShortestPathKernel,Mintz2009DistantSupervision}。这类方法在类别封闭、语料可控的场景下精度较高，且每一步的决策依据相对清晰。但问题也很直接：规则和特征的设计成本不低，且容易受限于特定领域的表达习惯；一旦换到新领域或遇到跨句关联，覆盖能力就会明显下降。更重要的是，实体识别和关系抽取如果采用流水线方式，前一步的错误会直接传导到后一步，整体召回和精度都容易受限。

深度学习的引入改变了这一局面。以 Lample 等\cite{Lample2016NeuralArchitecturesNER}和 Miwa 等\cite{Miwa2016EndToEndRE}的工作为代表，基于神经网络的方法开始直接从上下文中学习语义表示，减少了对人工词典和语言学特征的依赖。随后，统一结构生成模型进一步将抽取任务从分类问题转为文本到结构的预测，使实体识别与关系抽取得以在同一框架下联合完成\cite{Lu2022UIE,HuguetCabot2021REBEL}。相比特征工程阶段，这类方法在捕获复杂语境中的局部模式和语义关联上明显更具优势。不过，它们仍然依赖较为固定的 schema 设定和一定规模的标注数据。当实体类别体系、关系集合或文本风格发生变化时，往往需要重新训练或微调才能保持稳定。即便在统一生成框架下，长文本场景中的结构完整性、复杂约束条件下的边界识别等，也依然是容易出现偏差的环节。

近年来，大语言模型将知识抽取推进到一个新的阶段。一个核心变化在于，实体识别、关系抽取乃至更复杂的结构预测任务，都可以在自然语言指令的驱动下，以生成式的方式统一完成\cite{Xu2024LLMGenerativeIESurvey,Wang2023InstructUIE,DiazGarcia2025RESurveyLanguageModels}。这使得模型在处理长文本、跨句推理以及少样本迁移时展现出更强的适应性，也使得“从文本直接输出结构化结果”在技术上变得更为可行。但需要留意的是，大语言模型并没有自动解决领域知识抽取中的老问题。在低资源场景或面对未见过的类别时，性能仍然会明显下滑\cite{Deng2024LowResourceIESurvey}；与此同时，生成式输出的不确定性也在增加，冗余关系、实体边界漂移、类型匹配错误等现象并不少见。再加上大模型在高频批处理、本地化部署和长期运行中的资源消耗，现阶段的研究重点已不再只是追求模型的抽取能力，而是开始更多地关注训练数据的质量、输出结果的结构可控性，以及在真实工程场景中的可用性。

换句话说，面向领域知识图谱构建的知识抽取研究，正面临一个从“模型能不能抽出来”到“抽出来的结果能不能可靠入库”的转变。能否对抽取结果施加有效约束、能否保证其稳定性和可复用性，正在成为决定技术落地的关键。

\subsection{高质量可控知识抽取研究进展}

近年来，领域知识图谱的构建对知识抽取提出了新的要求，研究者开始将目光从单纯的单次抽取准确率，转向更全面的质量保障链条，这涉及训练数据的来源、结果的评估方式以及下游应用的有效利用。这一转变意味着，高质量的“银标”数据不再等同于标签数量的简单累积。例如，有学者将大语言模型在数据标注中的角色，定位在“生成-评估-筛选-利用”的整体流程中考察\cite{Tan2024LLMDataAnnotationSurvey}；而针对真实数据集的标注质量分析也表明，有效的质量控制依赖于一致性标准的设定、复核机制的建立以及验证流程的健全\cite{Klie2024AnalyzingAnnotationQualityManagement}。由此可见，面向领域的知识图谱构建，其数据基础必须同时满足准确性、完整性、一致性与可靠性这几个相互关联的条件。

为达成这一目标，现有研究从多个方向进行了探索。在人机协同方面，部分工作尝试将有限的人工精力优先投入到模型难以判断的“不确定样本”上，以期在标注成本与最终质量之间取得平衡\cite{Barbosa2025HILTS}。在模型能力提升上，有研究指出，为使通用模型在特定信息抽取任务上表现稳定，面向任务进行专项对齐往往必不可少\cite{Qi2024ADELIE}。此外，自验证方法在高专业文本抽取中展示了其价值：通过要求模型在输出结果的同时提供支撑证据并进行自我检查，能够在标注样本不足的情况下提升抽取的可靠性\cite{Gero2023SelfVerificationClinicalIE}。这些尝试共同指向一个核心认识：构建高质量训练数据是一个涉及“谁来标、标哪些、如何校”的系统性问题，而非单纯的“自动生成标签”技术。然而，目前的研究多集中于主动学习、人机协同或自验证等单个环节，尚缺乏一套能将提示优化、抽取生成、事实评估与错误回流整合为一个闭环，并直接服务于领域知识图谱构建需求的完整方案。可以说，如何为 Teacher 端系统性地生产高质量的银标数据，目前仍是一个有待解决的核心问题。

与数据质量问题并行的，是面向实际部署场景的模型轻量化研究。LoRA、QLoRA 等方法显著降低了模型适配领域时的显存占用和训练成本\cite{Hu2022LoRA,Dettmers2023QLoRA,Han2024PEFTSurvey}，使得在资源有限的环境下构建专用的知识抽取模型成为可能。同时，小语言模型的研究也表明，在需要本地部署、长期运行或应用于终端设备的场景中，小参数模型在响应延迟、资源开销和运维便捷性上具有天然优势\cite{Lu2024SmallLanguageModels}。这为 Student 端采用“小参数模型+参数高效微调”的路线提供了技术基础。但一个需要警惕的误区是：参数高效微调解决的是“如何低成本训练与部署”，它本身并不能保证模型输出的结果在语义和结构上是“可控”的。如果缺少对本体约束、关系合法性以及输出结果的选择机制，模型即使完成了抽取任务，其生成的内容也可能难以解析、无法满足图谱的模式要求，甚至无法稳定入库。

因此，对“可控抽取”的探讨，逐渐向如何将结构约束前置到生成过程的方向演进。有研究强调，知识图谱构建中的抽取系统应具备适应动态 schema 的能力，并在生成过程中显式考虑结构约束\cite{Ye2023SchemaAdaptableKGC}；另有研究尝试将领域本体知识引入抽取阶段，以从源头减少非法关系或不合逻辑三元组的产生\cite{vanCauter2024OntologyGuidedKGC}。在具体实现上，受限解码和语法约束生成技术进一步表明，与其依赖事后修补，不如在生成阶段就对输出结果的结构合法性进行控制\cite{Hemmer2023LazyKDecoding,Geng2023GrammarConstrainedDecoding}。此外，关于候选结果选择与规范化的研究指出，模型的单次生成结果并非最终可用的知识，对困难样本的重排序、结果的规范化处理，是连接模型输出与稳定入库的关键环节\cite{Ma2023LLMNotGoodFewShotIE,Zhang2024EDC}。相应地，面向文本到知识图谱生成的评价体系，也已将本体一致性、事实忠实性以及对幻觉的抑制，纳入“可控性”的考察范畴\cite{Mihindukulasooriya2023Text2KGBench}。不难看出，当前研究为可控抽取提供了多种可借鉴的技术手段，但多数工作仍各自聚焦于格式控制、一般 schema 感知或后验筛选，尚未将这些技术与小参数模型、参数高效微调以及具体的领域本体约束进行有机整合。

综合来看，现有研究至少存在三个需要进一步突破的缺口。其一，在数据构建端（Teacher 端），缺乏一个围绕领域知识图谱需求设计的、集提示优化、抽取生成、事实评估与错误回流于一体的闭环方法。其二，在模型部署端（Student 端），虽然证明了小参数模型结合微调的可行性，但“能部署”与“结果可控”之间仍存在断层，参数高效微调与结果结构控制尚未被统一设计。其三，现有的 schema 与本体约束研究，在支撑稳定入库方面仍有距离，尤其是在“实体类型-关系类型”的显式匹配、多候选结果的择优与规范化处理等环节，协同性不足。这些缺口的共同指向是：面向领域知识图谱构建的知识抽取，既需要在源头解决高质量银标数据的供给问题，也需要在中间环节建立起一套能兼顾参数高效微调、本体约束与多候选选择的协同可控机制。基于此，本文将从 Teacher 端（高质量数据构建）与 Student 端（可控抽取机制）两个层面展开后续研究。

\section{本文研究内容、创新点与技术路线}

针对领域知识图谱构建中的知识抽取问题，本文围绕一条贯穿“数据-抽取-应用”的主线展开：从低资源场景下的高质量训练数据构建，到面向本地化部署的参数高效抽取，再到满足本体约束的稳定入库与下游应用验证。具体而言，需要回应三个相互关联的工程性问题：当领域标注数据稀缺时，如何获得高置信度的训练样本并控制其质量；在算力和部署条件受限的情况下，能否以较低的参数量实现可用的领域抽取能力；抽取结果如何与领域本体对齐、形成合法三元组，从而支撑稳定入库，而不是仅停留在句子层面的“文本正确”。

围绕上述问题，本文的主要研究内容如下。
\begin{enumerate}
  \item 在数据构建层面（Teacher 端），以少量种子样本为起点，通过提示自优化、事实评估与回流纠错形成闭环机制，生成高置信、高一致性的银标数据。这一环节的目标是把数据质量控制从事后清洗前移到源头构建阶段。
  \item 在模型训练与推理层面（Student 端），首先以参数开源的小模型为基座，结合 QLoRA/PEFT 进行领域微调，在保留抽取能力的同时降低本地化部署与长期运行的门槛；随后在推理阶段引入本体约束与多候选选择机制，利用显式的实体类型与关系类型匹配规则抑制非法三元组，并通过多候选结果的综合评分提升输出一致性和稳定入库能力。
  \item 在应用验证层面，将上述抽取结果用于领域知识图谱的构建，并通过检索、简单问答等基础应用检验方法能否贯通从原始文本到结构化知识组织再到基础服务的完整流程。这里更关注的是方法在真实工程环节中的可落地性，而非单一评价指标上的增益。
\end{enumerate}

与现有研究相比，本文的差异主要体现在以下三个方面。
\begin{enumerate}
  \item 将高质量银标数据的构建从“模型训练后清洗”前移为“训练前的核心前提”，通过 Teacher 端的闭环机制将数据质量控制前置。
  \item 在 Student 端将参数高效微调与可控抽取统一在一个框架中，强调小参数模型、QLoRA、本体约束与多候选选择的协同设计，而非将模型压缩与控制策略视为两条独立的技术路径。
  \item 将知识图谱构建与下游应用验证纳入方法评估的必要环节，使研究目标从追求抽取精度扩展到支撑稳定入库与工程可用性，更贴近领域知识组织在实际部署中的需求。
\end{enumerate}

基于上述内容，本文的技术路线可以概括为：在 Teacher 端，依靠少量种子样本驱动提示自优化，结合事实评估与错误回流形成银标数据构建闭环；在 Student 端，以小参数模型为起点完成领域微调，并在推理阶段叠加本体约束与多候选选择，形成面向稳定入库的可控抽取机制；最后，将抽取结果用于领域知识图谱构建，通过基础检索与问答任务验证其工程可行性。

全文共分六章。第一章阐述研究背景与问题；第二章介绍相关技术与总体框架；第三章聚焦 Teacher 端银标数据的构建方法；第四章讨论 Student 端面向本体约束的参数高效可控抽取；第五章开展领域知识图谱构建与下游应用验证；第六章总结研究工作并讨论未来方向。

\section{本章小结}

本章从领域知识图谱构建对高质量可控知识抽取的实际需求出发，分析了现有研究在数据闭环构建、参数高效微调与可控抽取协同、以及面向稳定入库的一体化机制等方面的不足。在此基础上，明确了本文的整体研究思路：以 Teacher 端的银标数据构建和 Student 端的本体约束可控抽取为核心，以下游应用验证为支撑，形成贯通“数据-抽取-应用”的技术链条，为后续各章的具体工作奠定框架基础。
